\documentclass[a4paper]{article}

% Basic packages
\usepackage[fontset=fandol]{ctex}
\usepackage{amsmath, amssymb}
\usepackage{graphicx}
\usepackage[margin=2.5cm]{geometry}
\usepackage[most]{tcolorbox}
\usepackage{etoolbox}
\usepackage{listings}
\usepackage{hyperref}
\usepackage{booktabs}
\usepackage{subcaption}
\usepackage{float}
\usepackage{tikz}
\IfFileExists{pgfplots.sty}{
    \usepackage{pgfplots}
    \pgfplotsset{compat=1.18}
}{}

% Highlight boxes
\newtcolorbox{knowledgebox}[1]{
    enhanced, colback=blue!5!white, colframe=blue!75!black, colbacktitle=blue!75!black,
    coltitle=white, fonttitle=\bfseries, title=#1, attach boxed title to top left={yshift=-2mm, xshift=2mm},
    boxrule=1pt, sharp corners
}

\newtcolorbox{importantbox}[1]{
    enhanced, colback=yellow!10!white, colframe=yellow!80!black, colbacktitle=yellow!80!black,
    coltitle=black, fonttitle=\bfseries, title=#1, sharp corners
}

\newtcolorbox{warningbox}[1]{
    enhanced, colback=red!5!white, colframe=red!75!black, colbacktitle=red!75!black,
    coltitle=white, fonttitle=\bfseries, title=#1, sharp corners
}

% Listings
\lstset{
    language=Python,
    basicstyle=\ttfamily\small,
    keywordstyle=\color{blue},
    stringstyle=\color{red!60!black},
    commentstyle=\color{green!60!black},
    breaklines=true,
    frame=single,
    numbers=left,
    numberstyle=\tiny\color{gray},
    captionpos=b,
    extendedchars=false
}

% ============================================================
%  Lectio 讲义模板 —— 使用说明（给 LLM 看）
%  - 本模板是「参考骨架」，不是强制结构。根据输入材料的实际形态调整：
%      · slides → notes：可按幻灯片章节切 section，省略 \sourceduration。
%      · video  → notes：保留时长/链接，正文按时间线组织。
%      · text   → notes：来源元数据可留空或删除整个 titlepage。
%  - 必须保留：\usepackage 区块、tcolorbox 定义、compile.sh 编译方式。
%  - 可以改：\newcommand 默认值、titlepage 排版、section 结构、是否要 TOC。
%  - 无关字段请「删掉对应 \newcommand 与 titlepage 行」，而不是留 [在此填写...]。
% ============================================================

% Front-page metadata（用不到的字段可直接删除对应命令及 titlepage 中的引用）
\newcommand{\notetitle}{课程讲义}
\newcommand{\noteauthors}{Lectio}
\newcommand{\notedate}{\today}
\newcommand{\sourcechannel}{}
\newcommand{\sourcepublishdate}{}
\newcommand{\sourceduration}{}
\newcommand{\sourceurl}{}
\newcommand{\coverpath}{}

\begin{document}

\begin{titlepage}
\centering
{\Large 课程讲义\par}
\vspace{1.2cm}
{\huge\bfseries \notetitle\par}
\vspace{0.8cm}
{\large \noteauthors\par}
\vspace{0.3cm}
{\large \notedate\par}
\vspace{1.2cm}

\ifdefempty{\coverpath}{
{\small 请在生成讲义时填入封面图的本地路径。\par}
}{
\includegraphics[width=0.82\textwidth,height=0.45\textheight,keepaspectratio]{\coverpath}\par
}

\vfill
% 来源信息框：若对应字段为空会自动省略整行。
% 纯 slides/text 输入可以整段删除这个 tcolorbox。
\begin{tcolorbox}[width=0.9\textwidth, colback=black!2!white, colframe=black!60, sharp corners]
\ifdefempty{\sourcechannel}{}{\textbf{来源/作者}：\sourcechannel\par}
\ifdefempty{\sourcepublishdate}{}{\textbf{发布日期}：\sourcepublishdate\par}
\ifdefempty{\sourceduration}{}{\textbf{时长/篇幅}：\sourceduration\par}
\ifdefempty{\sourceurl}{}{\textbf{链接}：\href{\sourceurl}{\nolinkurl{\sourceurl}}\par}
\end{tcolorbox}
\end{titlepage}

% 短讲义（< 5 页）可以把下面两行删掉。
\tableofcontents
\newpage

%% --- 正文内容开始（以下 section 结构仅为示例，按实际材料组织即可） --- %%

% 约束（硬性）：
%   1) 图片放 \includegraphics，不要塞进 knowledgebox/importantbox/warningbox。
%   2) slides 来源的图在同页加脚注「来源：Slide P.N」；视频帧则「来源：mm:ss」。
%   3) 每个 \section 末尾放一个 importantbox 做费曼压缩（一句话版）。
%   4) 全文末尾 \section*{思考题} 放 3–5 个开放问题。
%
% 自由度：章节数量/标题/是否分小节、是否用 knowledgebox、是否加附录，按内容决定。

%% --- 正文内容结束 --- %%

\end{document}
